
\chapter{The eyes}

\textbf{Purpose:} Add occlusion along a ray to the reading map as a row type, measure whether it
reaches directions the harmonic rows cannot, and price the closure. \textbf{Scope:}
\texttt{beam\_rows.py}, \texttt{arm\_rank.py}, \texttt{reading\_rank.py},
\texttt{build\_sha\_clock\_view.py}.

A boundary reading to degree 8 over 256 sources has rank 81 and is blind in 175 directions. Adding
256 beams closes \textbf{all 175} and takes the reading to rank 256, the source count. The
closure is real, and not nominal: a random binary state recovers at \(6.5 \times 10^{-15}\) with a
condition number of 55.

\section{Why a beam is a different kind of row}

The viewer has carried these from the start and its own documentation states the mechanism: what
stops the carried beam is the bit being set, what casts the shadow is the stopping, and the pattern
on the wall is the state.

a beam is \emph{occlusion along a ray}, and therefore a line integral, where every instrument
already in the engine was a different shape.

\begin{center}
\begin{tabular}{@{}lll@{}}
\toprule
instrument & what it integrates & over what \\
\midrule
a harmonic row & the lit set against \(Y_{lm}\) & the whole sphere \\
an arm & the weight inside a region & a region \\
\textbf{a beam} & the weight lying along one line & \textbf{a line} \\
\bottomrule
\end{tabular}
\end{center}

A line integral is not in the span of region integrals. A beam is worth adding for that alone, and the claim
is a statement about linear algebra, and not about metaphor. The arms are touch; the beams are
eyes; the difference is which functional of the same state each one takes.

\section{What the beams buy}

\begin{center}
\begin{tabular}{@{}rrrrr@{}}
\toprule
beams & own rank & stacked rank & new directions & reach into kernel \\
\midrule
8 & 8 & 89 & 8 & 0.783320 \\
32 & 32 & 113 & 32 & 0.771013 \\
64 & 64 & 145 & 64 & 0.968161 \\
128 & 128 & 209 & 128 & \textbf{1.000000} \\
256 & 256 & \textbf{256} & \textbf{175} & 1.000000 \\
512 & 256 & 256 & 175 & 1.000000 \\
\bottomrule
\end{tabular}
\end{center}

\textbf{Every beam added is a new direction, with no waste at any count until the ceiling.} Eight buy
eight. A hundred and twenty eight buy a hundred and twenty eight. At 256 the stacked rank reaches the
source count and the harmonic kernel is gone.

Reach is the quantity to read: it is the largest principal angle cosine between the beam row space
and the harmonic kernel. A value of \(1.000000\) means some direction of the beam set lies entirely where the
harmonics are blind. It arrives at 128 beams and holds.

\textbf{The ceiling holds and is not a disappointment.} Five hundred and twelve beams still give 256,
because 256 is the number of sources. Beams cannot pass the dimension of the thing being read and are
not meant to; what they do is climb to it from a direction the harmonics do not have. That bound has
now been measured four independent ways: 1024 arms give 256, a precision sweep from 16 to 55 digits
gives 256, 512 beams give 256, and the kernel of the reading map empties at exactly 256.

\section{Both controls, and the negative one carries the weight}

\begin{center}
\begin{tabular}{@{}ll@{}}
\toprule
control & result \\
\midrule
the harmonic rows restacked on themselves & rank 81, \textbf{+0 new}, reach \(1.132 \times 10^{-15}\) into their own kernel \\
12 rows drawn \emph{from} the kernel & rank 93, \textbf{+12 new}, reach 1.000000 \\
\bottomrule
\end{tabular}
\end{center}

An instrument that reports a gain for a copy of its own rows reports a gain for anything. This one
reports zero for the copy and exactly twelve for twelve kernel rows. The count means what it says.

\section{Is the closure usable, which rank alone cannot answer}

This work has already been caught by the gap between completeness and usability: degree 15 is
complete at a least singular value of \(3.5 \times 10^{-3}\) and degree 16 is complete at
\(1.5 \times 10^{-1}\), a factor of forty four in noise tolerance between two maps of equal rank. The
closure therefore takes its grade from a round trip against a random \emph{binary} state, the form of a bit
pattern.

\begin{center}
\begin{tabular}{@{}lrrrr@{}}
\toprule
beam set & rank & least live & condition & recovery error \\
\midrule
harmonics alone & 81 & 4.365 & 1.058 & \(5.605 \times 10^{-1}\) \\
plus 256 random beams & 256 & \(1.387 \times 10^{-1}\) & \textbf{55} & \(\mathbf{6.516 \times 10^{-15}}\) \\
plus 256 fixed on bits & 256 & \(4.113 \times 10^{-2}\) & 138 & \(1.025 \times 10^{-14}\) \\
plus 256 fixed, signed & 256 & \(4.113 \times 10^{-2}\) & 138 & \(1.025 \times 10^{-14}\) \\
\bottomrule
\end{tabular}
\end{center}

The closure is real. A condition number of 55 is good and the state comes back at the arithmetic
floor.

\textbf{Fixing the beams on bits is worse than spreading them, which was not the expected direction.}
A beam aimed head on at its own source runs along that source's radial direction, a set of them is
more alike than a set aimed across the object, and alike rows are redundant rows. The least live value
falls from \(1.387 \times 10^{-1}\) to \(4.113 \times 10^{-2}\) and the condition number rises from 55
to 138. Both recover. The loss is in noise tolerance, and not in completeness. It is the same
lesson as the golden placement beating a random one on efficiency, read from the other side.

\textbf{Signed amplitudes change nothing, to every digit.} A linear span does not care about signs.
What a sign would change is a recovery constrained to non-negative combinations, which is a different
and harder problem than the least squares run here.

\section{A beam is a distribution}

One beam aimed at a single source touches \textbf{51 of 256} sources, with an entropy of 2.1581 nats,
about nine sources' worth of spread. a beam is not a probe of one bit even when it is aimed at one,
because it is a line integral.

\textbf{Calling it a wave function claims a phase it does not have.} A probability distribution is the
square of a wave function and the difference is phase. These rows are real and non-negative:
intensity, no phase, and two beams cannot interfere or cancel. The harmonic rows are the opposite,
complex coefficients whose phase is what torsion reads. So the engine holds both kinds, amplitude in
the harmonics and intensity in the beams, and which one is the right instrument depends on the object.
For a bit pattern, the beams. For a complex amplitude vector, the harmonics.

\section{Obliquity: the reading rises and the chord falls}

\begin{center}
\begin{tabular}{@{}rrrrr@{}}
\toprule
impact & chord & sources lit & row norm & obliquity \\
\midrule
0.0000 & 1.2400 & 17 & 1.2864 & \(0.0^\circ\) \\
0.3044 & 1.0803 & 17 & 1.3967 & \(29.4^\circ\) \\
0.4565 & 0.8390 & 28 & 1.6102 & \(47.4^\circ\) \\
0.6087 & 0.2354 & 20 & \textbf{1.9253} & \(79.1^\circ\) \\
0.6848 & 0.0000 & 14 & 0.9471 & \(90.0^\circ\) \\
\bottomrule
\end{tabular}
\end{center}

\textbf{The row norm rises about fifty percent as the beam goes oblique, then collapses once the ray
misses the shell.} The chord through the interior shrinks the whole way, as
\(2\sqrt{r^2 - b^2}\). The chord is therefore the wrong measure, and an argument from it gives the wrong answer.
A reading of this table from the chord contradicts its own rows.

\textbf{The sources sit on a shell, and not throughout a ball.} What a beam reads is time spent near the
shell. A grazing ray is nearly tangent to it, and runs alongside many sources at once. A central
ray crosses the shell twice, briefly, and spends the rest of its chord in an empty interior. Obliquity
buys contact even as it loses chord. That is a property of a shell placement and would not hold for
sources filling a volume.

The functional form is \textbf{not} named here. Over one decade of angle the rise is consistent with a
logarithm, with a reciprocal square root of the cosine, and with several other shapes. Naming it wants
a wider sweep or the wall reading, where an oblique beam also spreads its shadow over more area at one
over the cosine.

\section{What the model does not contain}

\textbf{Pure occlusion.} A source attenuates the ray passing it and nothing scatters. There is no
backscatter term, no second order path and no re-emission, any claim resting on backscatter is
outside what this measures. Adding it would be a further row type again, with its own controls.

\textbf{And the opacity is a free parameter, and not a derived one.} Attenuation goes as
\(\exp(-\tau \ell)\). At \(\tau\) near zero there is no signal; at \(\tau\) large every beam saturates
to ``blocked'' and carries exactly one bit, the silhouette. Information per beam peaks at intermediate
\(\tau\), around \(\tau \ell \approx 1\). Maximum opacity therefore \emph{minimizes} what a beam reports.
A \(\tau\) from a natural constant is reproducible, and not physical.

\section{What this does not establish}

It measures a property of the reading map, and not of any physical field. Nothing here shows that
information in any medium propagates along rays, and the beam is a model of occlusion chosen because
the viewer already drew it that way.

And rank is silent on worth. A large change in directions the map cannot resolve is still a large
number, and rank is a ceiling on what can be recovered, and not a statement about what is worth
recovering.
